\section{Forschungsagenda (1)}

Die moderne Film- und Medienbranche ist - national und international mit einem geschätzten Gesamtmarktwert im Jahr 2024 von 108.84 Milliarden USD \parencite{ZionMarketResearch} eine der wichtigsten und erfolgreichsten Industrien der heutigen Zeit.
Mit einer jährlichen Wachstumsrate von über 7\% ist es essentiell, das Sehverhalten der Kunden bei einem stetig wachsenden Wettbewerb zu analysieren.

Das Ziel dieses Projektes ist es, Nutzer, basierend nach ihrem Filmgeschmack, in gleichartige Gruppen zu unterteilen und diese anschließend bezüglich ausschlaggebender Faktoren für diese Einordnung zu untersuchen.
Zuerst werden Zuschauer mit gemeinsam gesehenen Filmen und angleichenden Wertungen gruppiert.
Diese Ergebnisse werden folgend mit den Genrevorlieben der Nutzer verglichen um mögliche Gemeinsamkeiten zwischen Filmgeschmack und Genre hervorzuheben.

Dabei steht die Frage im Vordergrund, ob rein datengetriebene Cluster mit den klassischen Kategorisierungen wie "Drama", "Action" oder "Horror" korrespondieren.

Es soll geprüft werden, ob bestimmte Nutzergruppen eine signifikante Neigung zu gewissen Genrekombinationen aufweisen, welche als Fundament dieser Gruppen fungieren.
Abschließend wird analysiert, inwieweit die zeitliche Dynamik das Sehverhalten beeinflusst und welche Arten von Filmen - Nischenfilme oder Blockbuster - für die Bildung solcher Gruppen von hoher relevanz sind.

In diesem Kontext wird untersucht, ob populäre Blockbuster aufgrund ihrer hohen Reichweite als "Rauschen" im Netzwerk fungieren, während spezielle Nischenfilme möglicherweise als stabilisierende Ankerpunkte dienen.

Der grundlegende Datensatz zur Auswertung der Forschungsfrage basiert auf einem Datenset von GroupLens.
Darin enthalten sind 10 Millionen Zuschauerbewertungen von 72.000 Nutzern auf 10.000 Filme, wodurch dieser Datensatz eine ausreichend hohe Dichte aufweist, um signifikante Muster erkennen zu können.

Zur Auswertung der Daten wird eine Netzwerkanalyse zwischen den Zuschauern basierend auf unterschiedlichen Featureschwerpunkten durchgeführt.
Nach der Bildung einer hinreichend befriedigenden Clusteranalyse folgt ein Vergleich der einzelnen Ergebnisse mit dem Ziel, eine Typologie des digitalen Filmgeschmacks zu entwerfen, die über rein algorithmische Empfehlungen hinausgeht und die dynamik des mondernen Medienkonsums wiederspiegelt.

